\chapter{Discussion} \label{sec:Discussion}

The system was evaluated as a proof of concept against the requirements of the problem statement (\Cref{sec:ProblemStatement}), verified through controlled bench and climate-chamber testing rather than full offshore qualification. This chapter interprets the results of \Cref{sec:Results} on that basis. It first assesses, sub-question by sub-question, what the work has established and what remains open. It then turns to the central tension running through the sensor development, between the sensitivity that makes the printed sensor usable at the small strains of interest and the reproducibility this same sensitivity costs, before examining how much confidence the test setup itself allows in the reported measurements. The chapter closes by considering what these results mean for turning the prototype into a deployable product, and which directions of future work would most effectively close the remaining gaps.

\section{Assessment Against the Research Question}
The first sub-question asked how 3D-printed sensor characteristics could be optimized for accurate strain transfer and sensitivity. The five-stage print optimization of \Cref{Baseline Single-Sensor Response} answers this concretely for the chosen conductive sensing material, PLA, raising the relative resistance change at \SI{800}{\micro\strain} from \num{0.08} to \num{0.38}, with the line depth setting responsible for the largest single gain. The resulting gauge factor of approximately \num{480}, and up to \num{750} for one arm of the assembled bridge, is more than two orders of magnitude above that of a metallic foil gauge~\cite{metal_gauge}, demonstrating that usable sensitivity at strains far below those previously reported for this sensor type~\cite{G.O.A.T} is achievable through print-process control alone, without a change of sensing principle.

The second sub-question, on integrating a strain measurement system within a $30\times\SI{30}{\milli\meter}$ metallic surface, is answered by the assembled module itself. The Wheatstone bridge pair (\SI{25}{\milli\meter} by \SI{6.25}{\milli\meter}, \Cref{Baseline Single-Sensor Response}), the PCB (\SI{20.83}{\milli\meter} by \SI{25.02}{\milli\meter}, \Cref{sec:final_pcb}), and the \textonehalf AA battery holder (\SI{34.4}{\milli\meter}~$\pm$~\SI{0.5}{\milli\meter} by \SI{17.5}{\milli\meter}~$\pm$~\SI{0.5}{\milli\meter}, \Cref{sec:battery}) together fit within the constraints of the nut surface. The drift characterization on a steel plate (\Cref{Material & Mounting-Surface Selection (Drift)}) further confirms that the sensor can be printed directly onto a metal surface, though this was not tested under strain. The serpentine spans only \SI{25}{\milli\meter} of the \SI{30}{\milli\meter} available on the nut face, leaving room to extend it for additional sensitivity, whereas the battery holder's \SI{34.4}{\milli\meter} dimension exceeds the nut face. If this overhang is not acceptable in the final mounting, a smaller holder must be found.

The third sub-question addressed wireless communication suited to acquiring data across the many bolted joints of a turbine. The Zigbee mesh network described in~\Cref{sec:zigbee_network} addresses the scale of the deployment directly, and the signal-strength results from~\Cref{sec:signal_strength} show that the encapsulated module sustains a link at ranges relevant to a turbine interior, reaching a maximum of \SI{20}{\meter} against \SI{25}{\meter} for the bare module (~\Cref{tab:DistanceRange}). The silicone potting therefore costs only about \SI{3}{\decibel} of range, recoverable with a single transmit-power step. Whether this link margin is sufficient through an entire turbine structure, however, was not tested in the field and remains unanswered.

The fourth sub-question asked how ultra-low-power electronics could support autonomous operation beyond ten years. The measured sleep- and run-mode currents translate, through the charge budget of ~\Cref{tab:charge_10y}, into an estimated ceiling of approximately 55 reports per day for a full decade within the derated capacity of the selected battery at the highest measured current draw (\SI{50}{\celsius}), with substantially more headroom at lower reporting rates or lower temperatures. This is a budget computed from measured currents and extrapolated over ten years, however, not a demonstrated ten-year lifetime.

The fifth sub-question, concerned with performance across operating temperature and long-term use, was answered through extensive characterization. The individual sensor and the assembled bridge were both cycled over the full \SIrange{-20}{50}{\celsius} range (\Cref{fig:TPUvsPLA_neg_deg}, \Cref{fig:FS_ADCogVout}), and the bridge was run continuously for approximately \SI{60}{\hour} under constant conditions (\Cref{fig:Slope1}) and under a smaller temperature excursion of \SIrange{20}{25}{\celsius} (\Cref{fig:sidsteTest}). These tests establish that the sensor and bridge remain functional and track applied strain throughout the full temperature range. It was not possible, however, to vary strain across the full operating range, so the response to temperature and to strain was characterized simultaneously only over the smaller \SIrange{20}{25}{\celsius} excursion. Combined with the test-setup limitations discussed in \Cref{sec:reliability}, which makes the underlying temperature and long-term drift of the sensor difficult to separate from confounding effects, this sub-question is only partially closed.

The final sub-question asked how encapsulation could resist water, oil, and grease while preserving strain transfer. A two-part silicone compound was selected and shown to preserve the sensor's strain response after potting (\Cref{sec:encapsulation}) while attenuating the wireless link only slightly (\Cref{sec:signal_strength}). Its resistance to water, oil, and grease ingress, and its durability over years of environmental cycling, was not tested directly, so this sub-question is answered at the level of feasibility rather than qualification.

Evaluated together, the sub-questions show that the core sensing and system-integration challenges posed by the research question have been demonstrated on the bench, while the questions of field performance, absolute long-term accuracy, and multi-year durability remain open, and are treated in turn in the remainder of this chapter.

\section{The Sensitivity–Reproducibility Trade-off} \label{sec:tradeoff}
The high sensitivity established above and the reproducibility limitations examined here share a common origin, and making it explicit accounts for the limitations of the sensor as much as for its central achievement. The sensitivity follows directly from operating the serpentine at the near-separation resting point of \Cref{Sensor Structure & the Small-Strain Operating Point}, where the asperity contacts between adjacent turns sit just short of losing contact and a strain of a few hundred microstrain still produces a measurable resistance change. The same proximity to separation leaves the baseline resistance of every printed element acutely sensitive to the print process, temperature, and time, which are precisely the influences the sensor must tolerate in service. This behavior is therefore best understood as a trade-off inherent to operating a constriction-resistance sensor at this point, rather than as a deficiency of the fabrication that a better printer would simply remove.

Three consequences follow. The first concerns the optimization itself, where the print-to-print variation in the asperity geometry is larger at this operating point than it would be for a sensor resting far from separation. The line depth, which fixes how far the nozzle is lowered toward the substrate and is the dominant setting in the five-stage sequence of \Cref{fig:opt_pla}, is among the hardest to hold constant between prints, since the nozzle-to-substrate spacing varies from one print to the next by amounts in the order of micrometers, and at the near-separation operating point a variation even this small shifts the resting contact enough to move the sensing characteristics noticeably. The practical difficulty is therefore not identifying which settings help, but reproducing the exact sensing characteristics from one print to the next, such that while the direction of each optimization is evident, the precise response of a given element is harder to guarantee.

The second consequence concerns the temperature compensation and the confidence in it. The Wheatstone half-bridge from \Cref{Wheatstone Half-Bridge Design} rejects the shared temperature response of its arms only to the extent that those arms are matched, and while two elements can in principle be printed to the same characteristics, the print-to-print variation present in this work makes such close matching difficult to achieve in practice. A bridge assembled from mismatched arms balances around its own offset rather than about zero, and a residual temperature sensitivity then survives compensation in proportion to that mismatch, as \Cref{fig:sidsteTest} shows directly, where the output envelope of the deliberately mismatched bridge tracks the ambient temperature rather than remaining flat. Improving the matching of the arms is therefore one of the clearest routes to a better-compensated bridge, and the difficulty is a recognized one for 3D printed Wheatstone bridges in general~\cite{3DprintWheatstoneStrain}, rather than particular to this sensor. The mismatch is not the whole story, however, since the fixture-induced strain discussed in \Cref{sec:reliability} couples into the bridge output during a thermal sweep and can't be separated from any temperature dependence of the sensor itself, the imperfect compensation observed here can then not be attributed to the arm matching alone.

The third consequence concerns the path to a product. As the process stands, each bridge carries its own offset and its own residual temperature coefficient, so a single design cannot be calibrated once and copied identically across units, and per-unit calibration or field commissioning becomes the more realistic manufacturing model in place of a uniform type-calibration. Per-unit calibration and the binning of sensors by measured characteristics are in any case established practice elsewhere in the sensor industry, so nothing in the printed sensor precludes either approach. This need is a consequence of the current print-to-print variation rather than a permanent feature, and if the printing was made reproducible enough for a sensor to emerge the same way each time, a single type-calibration applied once to the design could replace the per-unit step. How the printing might be brought to that point, and what it would mean for manufacture, is discussed in \Cref{sec:process_control}.

\section{Reliability and Validity of the Measurements} \label{sec:reliability}
The sensitivity of the sensors rests on a test bench built for this project rather than on calibrated commercial test equipment, and several aspects of that bench limit how far its results can be trusted as absolute measurements of sensor performance.

The actuation is open loop, and the applied displacement is computed from the number of microsteps commanded to the stepper driver, rather than measured independently by a displacement or force sensor at the point of contact with the DUT. The commanded travel was nonetheless checked manually against a mechanical micrometer, which confirmed that the stage delivers the intended displacement under normal operation, so backlash in the lead screw, friction in the linear stage, and compliance in the fixture do not compromise the reported values in ordinary room-temperature use. What this manual check cannot replace is a continuous, traceable readout at the point of contact, and these same effects become most pronounced under the temperature extremes of a thermal sweep, where expansion and contraction of the bench add to them. The nominal strain values reported throughout this thesis should therefore be read as commanded strain that has been confirmed against an external reference and found reliable under normal operation, while carrying a larger and less quantified uncertainty toward the ends of the temperature range.

The way each sensor is mounted into the actuator fixture (\Cref{Experimental Setup}) also affects its run-to-run reproducibility. 
A sensor is fastened into the fixture by hand and cannot be seated in precisely the same position, alignment, and pretension from one mounting to the next, so a remounted sensor presents a slightly different resting contact state to the actuator. 
This does not reflect a change in the sensor itself, which is unaffected by remounting, but a change in the mechanical state the mounting imposes on it. That state can be recovered by re-establishing the operating point through the stage, adjusting the applied pretension until the resting contact is restored, after which the sensor reproduces its earlier response. Reproducing a measurement across mountings therefore relies on this deliberate re-adjustment rather than on the mounting returning the sensor to the same state by itself, and the adjustment is in turn made less repeatable by the fixture being made of plastic, whose clamping interface has gradually worn under the repeated mounting and dismounting over the course of the project. 
Together with the break-in and contact-settling behavior described in \Cref{Sensor Structure & the Small-Strain Operating Point}, these effects mean that a sensor's earlier response can be reproduced after remounting, but only once the operating point has been re-established rather than automatically.

A further difference between the test bench and the deployed configuration lies in how strain is delivered to the sensor. On the bench the sensor is held between two clamps and strained by moving them relative to one another, so the sensing region spans a short unsupported length between the grips. In the intended application the sensor is instead bonded directly to the strained surface of the nut and is backed by that surface along its whole length, following the surface strain point for point. This distinction matters mainly under compression, where an unsupported span can bow out of plane, or buckle, rather than shortening uniformly, whereas a surface-bonded film is constrained by the substrate and cannot do so. The tensile straining used for the characterization in this work corresponds to the sensor's near-separation operating mode and is transferred cleanly through the clamps, so the response reported here is a valid characterization of that mode. Reproducing the compressive and relaxation behavior under a representative surface-bonded load is left to the improved test methodology discussed in \Cref{sec:improved_test}.

The test bench itself is not thermally inert. Over a full temperature sweep the lead screw, frame, and fixture expand and contract along with the device under test, and because the sensor is, by design, sensitive enough to register a few hundred microstrain, this bench-induced strain couples directly into the measured signal. This is the most likely explanation for the large, non-repeating excursions of the bridge output seen across the thermal cycle shown in \Cref{fig:FS_ADCogVout}, where the output at a given temperature differs by tens of millivolts between successive visits to the same setpoint, a pattern more consistent with mechanical strain from the expanding fixture, than with a fixed residual temperature dependence of the bridge. The same figure also shows abrupt spikes in the output that develop far more quickly than the slowly ramped chamber temperature; because any temperature-driven response, whether well compensated or not, can only track the temperature and must therefore vary as gradually as the temperature does, these rapid steps cannot originate in the compensation and instead indicate strain released suddenly as the fixture slips or settles. 
Because the sensor was not mechanically decoupled from this fixture-induced strain, the test can neither confirm nor rule out a temperature dependence of the bridge intrinsic to the sensor itself, as already noted in \Cref{sec:strain_sensing_performance}.

A related limitation is that strain could not be applied inside the climate chamber across the full range of temperature extremes explored. The effect of exposure from \SI{-40}{\celsius} to \SI{80}{\celsius} was therefore established by first exposing the sensor to a given temperature and only afterward straining it at room temperature (\Cref{fig:PLA_after_temp_limits}), rather than by straining it while held at that temperature. This shows that exposure alone does not degrade the sensor's subsequent response, but it does not demonstrate that the sensor performs identically while simultaneously strained and held at a temperature extreme, a gap that matters most at the coldest and hottest ends of the intended \SIrange{-20}{50}{\celsius} operating range.

None of these limitations were addressed within the scope of this thesis, since doing so would have required a calibrated load frame with independent displacement and force readout, together with a fixture capable of applying strain while thermally isolated from the surrounding environment, both beyond the equipment available this project. With such equipment, discussed further in \Cref{sec:improved_test}, the absolute accuracy of the reported strain values, and the separation of genuine sensor drift from rig and fixture effects, could be established with considerably more confidence.

\section{Implications for a Deployable Product}
Considered as a whole, the results of this chapter show the system to be a demonstrated proof of concept rather than a qualified product, and it is worth setting out clearly where that distinction currently appears.

The ten-year operating life, central to the research question, is established here as a measurement-based budget rather than a directly observed lifetime.
\Cref{sec:Current_Cons} shows that the measured currents are compatible with a decade of operation at realistic reporting rates, and because this projection rests on directly measured consumption rather than assumption, it forms a sound basis for the lifetime claim. 
A full ten-year run lies beyond the reach of any bench program, let alone a thesis, so confirming the budget in service is properly a task for long-term deployment. 
In the same way, the drift and temperature characterization from \Cref{sec:strain_sensing_performance} covers the hours to a few days timescales accessible on the bench and establishes the sensor's behavior over them, and the encapsulation has been shown to preserve sensing and RF performance immediately after potting (\Cref{sec:encapsulation}, \Cref{sec:signal_strength}). Extending these results to multi-year cycling and ingress is the validation that naturally follows rather than a gap in the present work. What will ultimately determine whether the system can detect a real loss of bolt preload in the field, is the absolute long-term strain accuracy and drift of the sensor, a quantity that lies beyond the reach of bench data of this duration and is therefore the central objective of the field validation to follow, for which the sensitivity and short-term stability established here provide the foundation.

Against this, the results also support the core premise that motivated the design. The unit cost of approximately \SI{80}{DKK} is low enough, relative to the value of monitoring a single high-consequence bolted joint, to make a one-module-per-bolt deployment plausible at the scale of a modern turbine's several thousand bolts~\cite{6000bolts}, in clear contrast to the proprietary, cost-constrained monitoring hardware discussed in \Cref{sec:previous_work}. Combined with the demonstrated sensitivity, wireless architecture, and low-power operation, this suggests that the remaining gaps are ones of qualification and field validation rather than of the underlying concept.